\chapter{Proposed Evaluation Methodology}
\label{Chapter4}

\section{Overview}

This thesis evaluates LLM-based Android malware detection methods under a
common framework: the same dataset(s), the same metrics
(Section~\ref{sec:metrics}), and a comparable evaluation protocol across
candidates, so that differences in reported performance reflect the method
rather than the evaluation setup. As of this writing, one candidate has a
working implementation: the LAMD-based method described in
Section~\ref{sec:lamd-candidate}, built on the LAMD framework introduced in
Chapter~\ref{Chapter2}. Additional candidates, drawn from the taxonomy of
LLM-based approaches identified in the Chapter~\ref{Chapter3} systematic
literature review (e.g. the RAG-based and fine-tuning-based families in
RQ1), are targeted for future integration into the same framework rather
than already implemented; this chapter is scoped to what has been built so
far and states explicitly where a candidate is still pending.

\section{Candidate Methods}

\textit{[TODO: Introduce the full set of candidate methods evaluated in this
thesis and how each was integrated into a common evaluation pipeline. One
candidate builds directly on the LAMD framework~\cite{Qian2025-ua}; its
development is detailed in Section~\ref{sec:lamd-candidate} below. Other
candidates (e.g. representative RAG- or fine-tuning-based systems from the
Chapter~\ref{Chapter3} taxonomy) should be described in their own subsections
here.]}

\subsection{Development of the LAMD-based Candidate}\label{sec:lamd-candidate}

This section reports the iterative engineering process behind the LAMD-based
candidate method, from a faithful baseline reimplementation through three
independent optimization attempts to the design currently used for
evaluation. Consistent with the methodology-focused scope of this chapter,
the discussion below concerns design rationale only; the corresponding
accuracy, cost, and latency measurements are reported in
Chapter~\ref{Chapter5}.

\subsubsection{Baseline: A Faithful LAMD Reimplementation}

No official implementation of LAMD~\cite{Qian2025-ua} is publicly available,
so a baseline was reconstructed directly from the paper's description of its
two-component design: \emph{key-context extraction}, which seeds analysis
from suspicious API calls and retains only the control- and data-dependent
code around them via backward program slicing, and \emph{tier-wise code
reasoning}, three LLM passes that build understanding bottom-up from
individual function summaries (Tier 1) to per-API intent (Tier 2) to a final
APK-level verdict (Tier 3), gated by a Factual Consistency Verification (FCV)
step between Tier 1 and Tier 2 that rejects function summaries whose Data
Relationship Coverage falls below the paper's threshold ($\theta = 0.95$).
This reimplementation is treated as the control condition throughout: every
sliced control-flow graph (CFG) is passed through all three tiers with no
caching and no early exit, exactly mirroring the paper's design without any
of the optimizations described below.

\subsubsection{Iterative Optimization Attempts}

Three independent directions were explored on top of the baseline, each
targeting the same underlying problem: every sliced CFG pays the full
Tier~1 + FCV + Tier~2 + Tier~3 cost regardless of how obvious the eventual
verdict is.

\textbf{Paper-fidelity corrections.} Before optimizing, two gaps against the
paper's own design were closed. First, the backward-slicing procedure now
recursively resolves parameters into caller methods when a variable is not
already defined within the current function, matching the paper's
requirement to recurse into callers until all variables are resolved.
Second, the FCV system prompt was corrected to reference the function under
verification explicitly, as the paper's own prompt specifies. These fixes
improve fidelity to the published method but do not by themselves change
cost or accuracy; they establish a more trustworthy baseline for the
optimizations that follow.

\textbf{Role-based early exit.} The first optimization classifies every
Tier-1 function summary with an additional \emph{role} label -- SOURCE,
SINK, BOTH, or NEUTRAL -- describing whether the function's suspicious API
usage looks like a data-collection point, a data-transmission point, both,
or neither. If every sliced CFG in an APK is classified NEUTRAL, the pipeline
exits early with a BENIGN verdict, skipping Tier~2 and Tier~3 entirely; if a
SOURCE-labeled CFG and a SINK-labeled CFG are both observed anywhere in the
APK, the pipeline exits early with a MALWARE verdict. This is a deliberately
minimal test of whether role-based shortcutting can reduce cost without
changing the verdict on clear-cut cases. It also has an obvious weakness:
a SOURCE and a SINK appearing \emph{anywhere} in the same APK are treated as
equally strong evidence regardless of whether they are actually related,
which risks a false positive whenever an application happens to contain two
unrelated suspicious-looking functions. This weakness directly motivated the
connection-verification design described next.

\textbf{Semantic caching with verified SOURCE/SINK connections.} The current
candidate design (Figure~\ref{fig:lamd-cache-pipeline}) keeps the role
classification idea but replaces the "any pair, anywhere" rule with an
explicit connection check, and adds a semantic cache to avoid repeating
Tier~1 and FCV work on near-duplicate code. A persistent cache stores each
CFG's verified Tier-1 summary alongside an embedding of the CFG text itself;
a new CFG whose embedding is sufficiently similar (cosine similarity
$\geq 0.88$) to a cached entry reuses the stored, already-verified summary
instead of issuing new Tier-1 and FCV calls. Once a SOURCE-labeled and a
SINK-labeled CFG have both been observed, the pipeline checks whether they
are actually connected before treating the pair as evidence, at one of three
escalating levels: (1) the two calls occur in the same method, which is
treated as strong structural evidence on its own, since the sink's backward
slice already includes the source's data path; (2) the two calls occur in
the same class, which is treated as suggestive but not conclusive on its
own; or (3) the two calls are reachable from one another in the
application's call graph within a bounded search radius. Levels 2 and 3
are additionally confirmed or rejected by a dedicated LLM verification call
before the pair is accepted as evidence; only level 1 is accepted without
further verification, since it is the only level with a structural
guarantee of a shared data path. If no connection is found, the candidate
pair is discarded and scanning continues rather than concluding MALWARE
prematurely.

This design was arrived at through repeated testing rather than settled in
advance. Two concrete refinements illustrate why the verification step
exists at all. First, an early version of the design treated same-class
co-location (level 2) as conclusive on its own, without the LLM verification
step -- on the reasoning that two suspicious methods sharing a class
plausibly share state. Testing against real-world samples showed this
assumption does not hold: two unrelated methods can share a utility class
with no data-flow relationship between them, and a same-class pair was
observed to trigger an early, incorrect benign classification before the
scan reached the application's actual malicious evidence. The design was
corrected so that level 2 requires the same LLM verification as level 3;
only level 1 retains its structural guarantee. Second, the suspicious-API
list used for key-context extraction was found to omit a legacy HTTP client
transmission API, which caused a small downloader-style sample to be
under-evidenced and misclassified; the API list was subsequently extended
with class-qualified matching so that the added API is recognized without
colliding with an unrelated, commonly used API of the same method name
elsewhere in the Android SDK. Both corrections are reported here as design
motivation; their measured effect on classification outcomes is reported in
Chapter~\ref{Chapter5}.

\textbf{Retrieval-augmented grounding (parallel exploration).} A fourth,
orthogonal direction grounds Tier~1 and Tier~3 prompts in a small,
hand-curated knowledge base of malware-relevant APIs and composite malware
profiles, retrieved via simple lexical (Jaccard) similarity rather than an
embedding-based retriever. Unlike the three optimizations above, this
direction does not aim to reduce cost; in preliminary testing it increased
both the number of sliced CFGs examined and the number of LLM calls issued,
while producing more concise and structured findings in the final report.
It is treated as a separate, not-yet-integrated exploration of explanation
quality rather than a candidate for the main cost/accuracy comparison in
Chapter~\ref{Chapter5}.

Table~\ref{tab:design-history} summarizes every design considered during
this iterative process, including intermediate variants superseded before
reaching the current design, for reference.

\begin{table}[htbp]
\centering
\caption{Every design considered for the LAMD-based candidate, including
superseded intermediate variants}
\label{tab:design-history}
\footnotesize
\begin{tabularx}{\textwidth}{lXXl}
\toprule
\# & Design & Outcome & Status \\
\midrule
1 & Baseline (faithful reimplementation) & Control condition; every CFG pays full Tier1+FCV+Tier2+Tier3 & Kept as baseline \\
2 & Role-based early exit (naive: any SOURCE + any SINK anywhere in the APK) & Coincidental unrelated pairs risk false positives & Superseded by connection-verification design \\
3a & Level-2 connection: shared Dalvik registers & Meaningless across method boundaries -- registers are not shared between functions & Discarded immediately \\
3b & Level-2 connection: same class, no LLM verification & Unsound -- caused a real false negative (unrelated methods sharing a generic utility class) & Fixed 2026-07-13 \\
3c & Level-2 connection: same class, with LLM verification (current) & Same treatment as Level 3 & Adopted \\
4a & Level-1 connection (same method): structurally guaranteed, no LLM verification & Still cost accuracy -- ``connected'' evidence is not necessarily malicious & Fixed 2026-07-18 \\
4b & Level-1 + uniform 0--10 LLM maliciousness score gate (threshold $\geq7$) (current) & +28\% calls on a 40-APK benchmark, accuracy 72.5\%$\to$80\% (this fix alone; row 8's fix accounts for the remaining 80\%$\to$85\%) & Adopted \\
5 & Semantic CFG cache (cosine similarity $\geq0.88$) & $-72\%$ to $-85\%$ calls/tokens (Chapter~\ref{Chapter5}) & Adopted \\
6 & Retrieval-augmented grounding (Jaccard similarity, no embeddings) & Increases calls rather than reducing them; more structured output & Parallel exploration, not integrated \\
7 & Suspicious-API list extensions (\texttt{HttpClient.execute()}, library allowlist rounds, OkHttp/Kawa/Cordova signature fixes) & Fixed several false negatives/positives from missing or over-broad API coverage & Adopted (some fixes not yet merged to the live batch runs) \\
8 & Tier-3 verdict-parsing fix (anchored \texttt{Verdict:} regex) & Fixed hedged-conclusion misparsing; 80\%$\to$85\% accuracy on the 40-APK benchmark & Adopted \\
\bottomrule
\end{tabularx}
\end{table}

\begin{figure}[htbp]
\centering
% Height-capped so the diagram plus its caption fit one page.
\adjustbox{max width=\textwidth,max totalheight=0.78\textheight}{
\begin{tikzpicture}[
  node distance=7mm and 10mm,
  every node/.style={font=\footnotesize},
  startstop/.style={rectangle, rounded corners, draw, fill=gray!10, minimum width=32mm, minimum height=7mm, align=center},
  process/.style={rectangle, draw, fill=blue!6, minimum width=32mm, minimum height=8mm, align=center},
  decision/.style={diamond, draw, fill=orange!8, aspect=2.2, inner sep=1pt, align=center},
  arr/.style={-{Latex[length=2mm]}}
]
  \node[startstop] (start) {For each sliced CFG,\\in order};
  \node[decision, below=of start] (cachehit) {Cache hit?\\($\geq 0.88$ sim.)};
  \node[process, left=14mm of cachehit, yshift=-9mm] (reuse) {Reuse cached\\summary (0 calls)};
  \node[process, right=14mm of cachehit, yshift=-9mm] (tier1) {Tier 1: summarize\\+ classify role};
  \node[process, below=of tier1] (fcv) {FCV: revise if\\DRC $< 0.95$};
  \node[process, below=13mm of cachehit] (role) {Track role: SOURCE /\\SINK / BOTH / NEUTRAL};
  \node[decision, below=of role] (pair) {Have a SOURCE\\\emph{and} a SINK item?};
  \node[process, left=20mm of pair] (next) {Next CFG};
  \node[decision, below=of pair] (conn) {Connection check\\(Level 1 / 2 / 3)};
  \node[process, below=9mm of conn] (verify) {Level 2/3: LLM verify\\CONFIRMED / BENIGN};
  \node[process, left=14mm of verify] (reset) {Reset pair,\\keep scanning};
  \node[process, right=16mm of verify, yshift=6mm] (strong) {Strong evidence\\-- break loop};
  \node[decision, below=14mm of verify] (allneutral) {All roles\\NEUTRAL?};
  \node[startstop, below left=10mm and -8mm of allneutral] (benign) {early\_neutral:\\BENIGN};
  \node[process, below right=10mm and -8mm of allneutral] (tier23) {Tier 2 + Tier 3 on\\non-NEUTRAL CFGs};
  \node[startstop, below=9mm of allneutral, yshift=-14mm] (final) {Final verdict};

  \draw[arr] (start) -- (cachehit);
  \draw[arr] (cachehit) -- node[above,xshift=-2mm]{yes} (reuse);
  \draw[arr] (cachehit) -- node[above,xshift=2mm]{no} (tier1);
  \draw[arr] (tier1) -- (fcv);
  \draw[arr] (fcv.south) .. controls +(0,-6mm) and +(18mm,6mm) .. (role.east);
  \draw[arr] (reuse.south) .. controls +(0,-6mm) and +(-18mm,6mm) .. (role.west);
  \draw[arr] (role) -- (pair);
  \draw[arr] (pair) -- node[above]{no} (next);
  \draw[arr] (next.north) .. controls +(0,20mm) and +(-30mm,0) .. (cachehit.west);
  \draw[arr] (pair) -- node[right]{yes} (conn);
  \draw[arr] (conn) -| node[near start, above]{none} (reset.north);
  \draw[arr] (conn) -- node[right]{Level 2/3} (verify);
  \draw[arr] (conn) -| node[near start, above]{Level 1} (strong.north);
  \draw[arr] (verify) -- node[pos=0.3, above, fill=white, inner sep=1pt]{CONFIRMED} (strong.south);
  \draw[arr] (verify) -- node[above]{BENIGN} (reset);
  \draw[arr] (reset.west) .. controls +(-10mm,0) and +(-10mm,0) .. (next.south);
  \draw[arr] (strong.south) |- (allneutral.east);
  \draw[arr] (verify.south) |- ++(0,-8mm) -| (allneutral);
  \draw[arr] (allneutral) -- node[above,xshift=-2mm]{yes} (benign);
  \draw[arr] (allneutral) -- node[above,xshift=2mm]{no} (tier23);
  \draw[arr] (benign.south) .. controls +(0,-6mm) and +(-14mm,6mm) .. (final.west);
  \draw[arr] (tier23.south) .. controls +(0,-6mm) and +(14mm,6mm) .. (final.east);
\end{tikzpicture}
}
\caption{End-to-end flow of the current LAMD-based candidate
(cache-and-verified-connection design). Each sliced CFG either hits the
semantic cache or goes through Tier~1 and FCV; once a SOURCE and a SINK item
have both been seen, a three-level connection check decides whether to treat
them as evidence before the pipeline falls through to Tier~2/3 or an early
verdict.}
\label{fig:lamd-cache-pipeline}
\end{figure}

\section{Datasets}

Development and debugging of the LAMD-based candidate used a small,
purpose-built sample of 12 APKs downloaded live via the AndroZoo API: 7
malware samples (\texttt{vt\_detection} $\geq 15$) and 5 benign samples
(\texttt{vt\_detection} $= 0$, sourced from the Google Play market), each
between 200KB and 1.3MB to keep pipeline runtime tractable during
iteration. This set is a development/debugging sample used to validate and
compare pipeline behavior across design changes (Chapter~\ref{Chapter5}); it
is distinct from the original LAMD paper's own evaluation set, which spans
13,794 AndroZoo samples labeled through VirusTotal ($>4$ engine detections)
across three temporally ordered test periods (2020--2023) to measure
robustness to distribution drift~\cite{Qian2025-ua}. A larger-scale
evaluation dataset for the final cross-candidate comparison has not yet
been finalized. Separately, an unfiltered random sample of 750 AndroZoo
APKs (no confidence threshold, no class balancing) is being run against two
additional self-hosted backends as an in-progress robustness check on the
candidate's false-positive behavior outside the curated 12-APK set; this is
exploratory and distinct from the finalized cross-candidate dataset above --
see Section~\ref{sec:threats-ch5}.

\section{Evaluation Metrics}\label{sec:metrics}

Detection quality is measured with Accuracy, Precision, Recall, and
F1-score, plus False Positive Rate (FPR) and False Negative Rate (FNR) to
separate the two error modes given the class imbalance typical of malware
datasets. Efficiency is measured separately via LLM call count, token usage,
wall-clock runtime per APK, and cost per APK in USD, since LLM-based methods
trade detection quality against non-trivial inference cost
(Chapter~\ref{Chapter3}, RQ2).

\section{Experimental Setup}

\textit{[TODO: Describe prompts/configurations and the procedure used to run
each candidate method once additional candidates are integrated.]}

The LAMD-based candidate is implemented in Python, using androguard for
static analysis and backward slicing, and a persistent semantic cache
(ChromaDB with its default ONNX embedding model) for the CFG-level cache
described in Section~\ref{sec:lamd-candidate}.

\textbf{LLM backend note.} The LAMD paper's own experiments use
GPT-4o-mini~\cite{Qian2025-ua}. For this thesis, all LAMD-based candidate
runs instead use a different backend model, accessed through a CLI-driven
subprocess call rather than a directly billed API, after the originally
planned API-based backend became impractical to sustain during development.
This substitution should be read as a methodology caveat: results produced
in Chapter~\ref{Chapter5} are not directly comparable to the original
paper's published figures on model choice alone. It is, however, consistent
with the original paper's own discussion of API cost as a practical
deployment constraint and its suggestion that cheaper or locally hosted
models are a natural next step.
